% Pursuing Stacks (cotes 134-2 to 134-8): archive cote 134-5, site page 21 (PDF page 22), typescript page 208.
% typescript. Draft compared: ps4.tex:950-1010.
% AI-transcribed by Claude Opus 5.5 (claude-opus-5-5, adaptive thinking, effort high), Message Batches API, batch ps2 (msgbatch_01NVwr1Yuo8qJf5rykeGBXPK), 2026-11-12; not yet checked by eye.
\scanpage{134-5/21}
\mspage{208}

ting to any topological group it's``\uline{classifying space}'',  defined ``up to homo-%
topy''.  I suspect the same should hold in any elementary modelizer $A^{\wedge}$,
$A$ a test category,  at least in the ``strict'' case\inserted{namely,} when $A^{\wedge}$\hand{/}is totally asphe-%
ric,  i.e.  the canonical functor  $A^{\wedge}\hand{$\longrightarrow$}$ (Hot)  commutes to finite products.
The corresponding statement for the basic localizer (Cat) itself however is
definitely false\unsure{.|,}  Group objects in (Cat) are indeed very interesting and
well-known beings  (introduced,  I understand from Ronnie Brown, by Henry
Whitehead long time ago,  under the somewhat misleading name of ``crossed
modules''),  yet they embody not arbitrary pointed 0-connected homotopy
types,  but merely those for which \hand{$\pi$}$_i(X) = 0$ for $i \hand{$>$} 2$ .  Thus we get only
2-truncated homotopy types  -  and presumably,  \xout{xxxxx} starting with n-Cat
instead of (Cat) as a modelizer,  we then should get \unsure{(n+1)}-truncated homoto-%
py types.  This ties in with the observation that taking group objects\inserted{either} in
(Cat),  or in the full subcategory  (Groupoids)  of the latter,  amounts to
the same  -  and similarly surely for (n-Cat);  on the other hand it has
been kind of clear from the very beginning of this reflection that \xout{xxxxxx}\inserted[typed]{at any}
\uline{finite} level,  groupoids and n-groupoids (n finite) will only yield \uline{trunca-}%
\uline{ted} homotopy types.

A related intriguing question is when exactly does a modelizer (M,W)
give rise to a \uline{Dold-Puppe theorem}  -  namely when do we get an actual \uline{equi-}%
\uline{valence} between the category of \uline{abelian} group objects of M,  and the catego-%
ry of chain complexes of abelian groups ?  \over{t}{T}he original statement was in
case  $M = \hand{$\Delta$}^{\wedge} =$  semisimplicial complexes,  and doubtlessly it was one main
impetus for the sudden invasion of homotopy and cohomology by semisimplicial
calculus  -  so much so it seems that for many people,  ``homotopy'' has be-%
come synonymous to ``semisimplicial algebra''.  The impression that semi-sim-%
plicial complexes is the \unsure{God}-given ground for doing homotopy and even cohomo-%
logy,  comes out rather strong also in Quillen's foundational notes, and
\inserted{in} Illusie\hand{'}s thesis.  Still there are too some cubical theory chaps I heard,
who surely must have noticed long ago that the Dold-Puppe theorem is valid
equally \xout{inxthe} for cubical complexes (I could hardly imagine that it\inserted{possibly} \hand{could}'
\xout{notxbe} nt).  Now it turns out that semisimplicial and cubical complexes are
part of a trilogy,  together with\inserted{so-called} ``\uline{hemispherical} \uline{complexes}'' ,  which look
a lot simpler still,  with just two boundary operators and one degeneracy
in each dimension.  They can be viewed as  embodying the ``primitive struc-%
ture'' of an \hand{$\infty$}- groupoid,  the boundary operators being the``source'' and
``target'' maps,  and the degeneracy the map associating to any i-object the
corresponding ``identity''.  \hand{(} I hit upon this structure among the first
examples of test categories and elementary modelizers,  and have been told
since by Ronnie Brown that he has\inserted{already} \hand{known} for a while these models,  under

\margin{21}
